\documentclass[12pt,a4paper]{article}

% --- Idioma y codificación ---
\usepackage[spanish]{babel}
\usepackage[utf8]{inputenc}
\usepackage[T1]{fontenc}
\usepackage{titlesec}
\titleformat{\section}
  {\normalfont\large\bfseries}
  {\thesection}{1em}{}
\titleformat{\subsection}
  {\normalfont\normalsize\bfseries}
  {\thesubsection}{1em}{}

% --- Márgenes ---
\usepackage[a4paper,
  top=2.5cm,bottom=2.5cm,
  left=2.2cm,right=2.2cm,
  columnsep=0.7cm
]{geometry}

% --- Gráficos, tablas, etc. ---
\usepackage{graphicx}
\usepackage{subfig}
\usepackage{physics}
\usepackage{csquotes}
\usepackage{float}
\usepackage{amsmath}
\usepackage{amssymb}
\usepackage{booktabs}
\usepackage{siunitx}
\usepackage{array}

% --- Dos columnas tipo paper ---
\usepackage{multicol}

% --- Índice y enlaces ---
\usepackage[hidelinks]{hyperref}
\addto\captionsspanish{\renewcommand{\contentsname}{Índice}}

% --- Encabezados y pies ---
\usepackage{fancyhdr}
\pagestyle{fancy}
\fancyhf{}
\lhead{Red de Difracción}
\rhead{\leftmark}
\cfoot{\thepage}
\setlength{\headheight}{14.62pt}
\addtolength{\topmargin}{-0.12pt}

% --- Bibliografía ---
\usepackage[backend=biber,style=ieee]{biblatex}
\addbibresource{bibliografia.bib}

\let\cleardoublepage\clearpage
\raggedcolumns

\begin{document}

% =========================
% Portada (1 columna)
% =========================
\begin{titlepage}
  \centering
  \vspace*{2cm}
  {\LARGE \textbf{Red de Difracción}\par}
  \vspace{1cm}
  {\large Asignatura: Laboratorio de Óptica\par}
  \vspace{0.5cm}
  {\large Autor: Luis López Nasser\par}
  {\large Fecha: 21/04/2026\par}
  \vfill
  {\large Universidad Unie\par}
\end{titlepage}

% =========================
% Índice (1 columna)
% =========================
\tableofcontents
\clearpage

% =========================
% Contenido (2 columnas)
% =========================
\begin{multicols}{2}

\section{Introducción}

\subsection{Marco Teórico}

Una \textbf{red de difracción} es un elemento óptico formado por una placa con un gran número de líneas paralelas grabadas a distancia constante $d$, denominada periodo o constante de la red. Cuando un haz monocromático de longitud de onda $\lambda$ incide sobre la red, la luz se difracta e interfiere produciendo un patrón de máximos brillantes a ángulos bien definidos.

\bigskip
La distribución de intensidad resulta de dos contribuciones: la difracción de cada rendija individual (de anchura $b$) y la interferencia entre las $N$ rendijas:

\begin{equation}
  I = I_0
  \left[\frac{\sin(\pi b\sin\theta/\lambda)}{\pi b\sin\theta/\lambda}\right]^2
  \left[\frac{\sin(\pi N d\sin\theta/\lambda)}{\sin(\pi d\sin\theta/\lambda)}\right]^2.
  \label{eq:intensidad}
\end{equation}

Los máximos principales ocurren cuando el denominador del segundo factor se anula, lo que conduce a la \textbf{ecuación de la red} para incidencia oblicua con ángulo $\varepsilon$:

\begin{equation}
  d(\sin\theta - \sin\varepsilon) = m\lambda, \quad m = 0,\pm1,\pm2,\ldots
  \label{eq:red_general}
\end{equation}

Para incidencia normal ($\varepsilon = 0$) se simplifica a

\begin{equation}
  d\sin\theta = m\lambda, \quad m = 0,\pm1,\pm2,\ldots
  \label{eq:red_normal}
\end{equation}

donde $m$ es el orden de difracción y $\theta$ el ángulo de observación respecto a la normal.

\bigskip
En la aproximación paraxial ($\theta$ pequeño), $\sin\theta \approx \tan\theta = x_m/z$, donde $x_m$ es la posición del máximo de orden $m$ en la pantalla y $z$ la distancia red-pantalla. Entonces:

\begin{equation}
  x_m = \frac{m\lambda z}{d},
  \label{eq:posicion}
\end{equation}

de modo que la separación entre el orden $+m$ y el orden $-m$ es

\begin{equation}
  \Delta x_m = x_m - x_{-m} = \frac{2m\lambda z}{d}.
  \label{eq:separacion}
\end{equation}

Representando $\Delta x_m$ frente a $m$ se obtiene una recta de pendiente $2\lambda z/d$, que permite determinar el periodo $d$ de la red.

\bigskip
Cuando la luz incidente contiene varias longitudes de onda, para cada orden $m$ cada componente produce su máximo a un ángulo diferente $\theta(\lambda)$, lo que permite separar el espectro. Para un orden dado, a mayor $\lambda$ mayor ángulo de desviación.

\subsubsection{Goniómetro}

El goniómetro consta de un brazo fijo (fuente, rendija, lente colimadora y soporte para la red) y un brazo móvil con anteojo (lente objetivo, retículo en cruz y ocular). La red se coloca en el centro del instrumento. Al no conocerse con precisión el ángulo de incidencia $\varepsilon$, se miden los ángulos $\theta_{+m}$ y $\theta_{-m}$ a izquierda y derecha del orden central, y el ángulo de difracción se obtiene como

\begin{equation}
  \theta = \frac{|\theta_{+m} - \theta_{-m}|}{2},
  \label{eq:goniometro}
\end{equation}

eliminando así la incertidumbre sobre $\varepsilon$. Con este $\theta$ y la ecuación~(\ref{eq:red_normal}) se calcula directamente $\lambda$.

\subsection{Objetivos}

\begin{itemize}
  \item Caracterizar una red de difracción determinando su periodo $d$ y el número de líneas por mm mediante el patrón láser para dos distancias.
  \item Determinar la longitud de onda de las líneas espectrales de la lámpara de sodio usando el goniómetro con una red de 300 líneas/mm.
  \item Comparar los valores obtenidos con los teóricos y calcular el error relativo.
\end{itemize}

\section{Materiales}

\begin{itemize}
  \item Fuente de luz láser ($\lambda$ conocida).
  \item Dos redes de difracción a caracterizar.
  \item Lámpara de sodio.
  \item Goniómetro (brazo fijo + brazo móvil con anteojo).
  \item Pantalla y regla milimétrica.
\end{itemize}

\section{Procedimiento}

\subsection{Caracterización de la red con láser}

\begin{enumerate}
  \item Colocar la red perpendicularmente al haz láser y observar varios órdenes de difracción en la pantalla.
  \item Repetir con la segunda red; seleccionar la que permita medir más órdenes con precisión.
  \item Para dos distancias $z$ distintas, medir $\Delta x_m$ para los órdenes $m = \pm1, \pm2, \pm3, \pm4$.
  \item Representar $\Delta x_m$ frente a $m$ y ajustar linealmente para obtener $d$ a partir de la pendiente.
\end{enumerate}

\subsection{Medida con el goniómetro (lámpara de sodio)}

\begin{enumerate}
  \item Colocar la red de 300 líneas/mm en el goniómetro.
  \item Registrar el ángulo $\theta_0$ con el anteojo centrado en el orden cero.
  \item Para cada línea espectral y para los órdenes $m = +1, +2, +3$ y $m = -1, -2, -3$, centrar la línea en el retículo y anotar $\theta_{\pm m}$.
  \item Calcular $\theta = |\theta_{+m} - \theta_{-m}|/2$ y obtener $\lambda$ con la ecuación~(\ref{eq:red_normal}).
\end{enumerate}

\section{Análisis y Discusión}

\subsection{Caracterización de la red: periodo $d$}

Usando la ecuación~(\ref{eq:separacion}), la pendiente del ajuste lineal de $\Delta x_m$ frente a $m$ es $p = 2\lambda z / d$, de donde

\begin{equation}
  d = \frac{2\lambda z}{p}, \qquad N = \frac{1}{d}.
  \label{eq:d_ajuste}
\end{equation}

% TODO: rellenar tablas con datos experimentales.

\subsubsection*{\texorpdfstring{Distancia $z_1$}{Distancia z1}}

\begin{table}[H]
  \centering
  \caption{Separación $\Delta x_m$ para $z_1$.}
  \label{tab:red_z1}
  \small
  \begin{tabular}{ccccc}
    \toprule
    $m$ & $\pm1$ & $\pm2$ & $\pm3$ & $\pm4$ \\
    \midrule
    $\Delta x_m$ & & & & \\
    \bottomrule
  \end{tabular}
\end{table}

\begin{equation*}
  z_1 \pm \Delta z_1 = \qquad d_1 \pm \Delta d_1 =
\end{equation*}

\subsubsection*{\texorpdfstring{Distancia $z_2$}{Distancia z2}}

\begin{table}[H]
  \centering
  \caption{Separación $\Delta x_m$ para $z_2$.}
  \label{tab:red_z2}
  \small
  \begin{tabular}{ccccc}
    \toprule
    $m$ & $\pm1$ & $\pm2$ & $\pm3$ & $\pm4$ \\
    \midrule
    $\Delta x_m$ & & & & \\
    \bottomrule
  \end{tabular}
\end{table}

\begin{equation*}
  z_2 \pm \Delta z_2 = \qquad d_2 \pm \Delta d_2 =
\end{equation*}

% TODO: insertar gráfica del ajuste lineal para las dos distancias.

\begin{figure}[H]
  \centering
  % \includegraphics[width=1\columnwidth]{imagenes/ajuste_red}
  \caption{Ajuste lineal de $\Delta x_m$ frente a $m$ para las dos distancias.}
  \label{fig:ajuste_red}
\end{figure}

\begin{table}[H]
  \centering
  \caption{Resultados de $d$ y $N$ para cada distancia.}
  \label{tab:resultados_red}
  \small
  \begin{tabular}{cccc}
    \toprule
    $z$ & Pendiente $p$ & $d \pm \Delta d$ & $N$ (l/mm) \\
    \midrule
    $z_1$ & & & \\
    $z_2$ & & & \\
    \bottomrule
  \end{tabular}
\end{table}

% TODO: desarrollar el cálculo de errores por propagación.

\subsection{\texorpdfstring{Longitud de onda de la lámpara de sodio}{Longitud de onda de la lampara de sodio}}

Se usa la red de 300 líneas/mm ($d = 1/300\ \text{mm}$). El ángulo de orden cero es $\theta_0 \pm \Delta\theta_0$.

% TODO: rellenar con datos del goniómetro.

\begin{table}[H]
  \centering
  \caption{Ángulos medidos con el goniómetro (sodio).}
  \label{tab:goniometro}
  \small
  \begin{tabular}{cccc}
    \toprule
     & $m=+1$ & $m=+2$ & $m=+3$ \\
    \midrule
    $\theta_{+m}$ & & & \\
    \midrule
     & $m=-1$ & $m=-2$ & $m=-3$ \\
    \midrule
    $\theta_{-m}$ & & & \\
    \bottomrule
  \end{tabular}
\end{table}

Para cada orden se calcula $\theta = |\theta_{+m} - \theta_{-m}|/2$ y se aplica la ecuación~(\ref{eq:red_normal}):

\begin{equation}
  \lambda = \frac{d\sin\theta}{m}.
  \label{eq:lambda}
\end{equation}

\begin{table}[H]
  \centering
  \caption{Longitudes de onda calculadas por orden.}
  \label{tab:lambda}
  \small
  \begin{tabular}{cccc}
    \toprule
    Orden $m$ & $\theta$ & $\lambda$ & $\Delta\lambda$ \\
    \midrule
    1 & & & \\
    2 & & & \\
    3 & & & \\
    \bottomrule
  \end{tabular}
\end{table}

\begin{equation*}
  \lambda \pm \Delta\lambda = \quad \text{(valor ponderado)}
\end{equation*}

% TODO: desarrollar el cálculo de error propagado de lambda.

\subsection{Comparación con valores teóricos}

% TODO: anotar los valores nominales de d y de lambda del sodio.

\begin{table}[H]
  \centering
  \caption{Comparación experimental vs.\ teórico.}
  \label{tab:comparacion}
  \small
  \begin{tabular}{lccc}
    \toprule
    Magnitud & Exp. & Teórico & Error rel. \\
    \midrule
    $d$ (mm) & & & \\
    $\lambda_{\text{Na}}$ (nm) & & 589{,}3 & \\
    \bottomrule
  \end{tabular}
\end{table}

% TODO: comentar si los resultados son razonables y discutir fuentes de error.

El doblete del sodio tiene sus dos líneas a \SI{589.0}{\nano\meter} y \SI{589.6}{\nano\meter}; si el goniómetro no tiene resolución suficiente para separarlas, se mide una longitud de onda efectiva próxima a \SI{589.3}{\nano\meter}.

\section{Conclusión}

% TODO: redactar conclusiones con los valores numéricos obtenidos.

\begin{enumerate}
  \item Se ha caracterizado la red de difracción mediante el ajuste lineal de $\Delta x_m$ frente a $m$ para dos distancias distintas, obteniendo un periodo $d \pm \Delta d$ consistente en ambos casos y un número de líneas $N = 1/d$.

  \item Con el goniómetro y la red de 300 líneas/mm se ha medido el ángulo de difracción de las líneas espectrales del sodio para los órdenes $m = 1, 2, 3$, obteniendo una longitud de onda ponderada $\lambda \pm \Delta\lambda$ compatible con el valor teórico de \SI{589.3}{\nano\meter} dentro de la incertidumbre experimental.

  \item Las principales fuentes de error son la resolución de la escala angular del goniómetro, la incertidumbre en la lectura de posiciones en la pantalla y la posible no-perpendicularidad del haz sobre la red.
\end{enumerate}

\end{multicols}

\printbibliography

\end{document}
